\section{مدول‌ها}

% Source: sources/lectures/04 Module.pdf | PDF page: 1 | item: left and right module definition
\begin{frame}{مدول چپ و مدول راست}
    \label{lec04:module-definition}
    \begin{definition}[مدول چپ]
        فرض کنید $R$ حلقه‌ای دلخواه و $(M,+)$ گروهی آبلی باشد.
        ضرب اسکالر $R\times M\to M$ را با $(r,m)\mapsto rm$ نشان می‌دهیم.
        برای هر $r,s\in R$ و $m,m_1,m_2\in M$ فرض می‌کنیم:
        \[
        \begin{aligned}
            r(m_1+m_2)&=rm_1+rm_2,\\
            (r+s)m&=rm+sm,\\
            (rs)m&=r(sm).
        \end{aligned}
        \]
        در این صورت $M$ یک \textbf{$R$-مدول چپ} است.
        مدول راست نیز به همین نحو تعریف می‌شود.
    \end{definition}
\end{frame}

% Source: sources/lectures/04 Module.pdf | PDF page: 1 | item: unitary module and first two examples
\begin{frame}{مدول یکانی و نمونه‌های نخست}
    \label{lec04:unitary-module}
    \begin{definition}[مدول چپ یکانی]
        اگر $R$ حلقه‌ای یکدار و $M$ یک $R$-مدول چپ باشد و
        \[
            1m=m\qquad(\forall m\in M),
        \]
        آنگاه $M$ را \textbf{یکانی} می‌نامیم.
    \end{definition}
    \begin{example}
        هر فضای برداری بر میدان $F$ یک $F$-مدول چپ یکانی است.
    \end{example}
    \begin{example}
        هر حلقهٔ $R$ با عمل $r\cdot s=rs$ یک $R$-مدول چپ است.
    \end{example}
\end{frame}

% Source: sources/lectures/04 Module.pdf | PDF page: 1-2 | item: abelian groups as Z-modules
\begin{frame}{گروه‌های آبلی به‌عنوان مدول}
    \label{lec04:abelian-z-module}
    \begin{example}
        هر گروه آبلی $(A,+)$ یک $\mathbb Z$-مدول چپ است، با عمل
        \[
            n\cdot a=
            \begin{cases}
                \underbrace{a+\cdots+a}_{n\text{ بار}} & n>0,\\[3pt]
                0 & n=0,\\[3pt]
                -\underbrace{(a+\cdots+a)}_{-n\text{ بار}} & n<0.
            \end{cases}
        \]
    \end{example}
    \begin{example}
        هر گروه آبلی $(M,+)$ با تعریف زیر یک $R$-مدول می‌شود:
        \[
            rm=0\qquad(\forall r\in R,\ \forall m\in M).
        \]
    \end{example}
\end{frame}

% Source: sources/lectures/04 Module.pdf | PDF page: 2 | item: further examples and the Z over R question
\begin{frame}{نمونه‌ها و یک پرسش}
    \label{lec04:scalar-examples}
    \begin{example}
        $\mathbb R[x]$ با ضرب معمول اسکالر یک $\mathbb R$-مدول چپ است.
    \end{example}
    \begin{example}
        $\mathbb R$ با ضرب معمول اعداد صحیح یک $\mathbb Z$-مدول است.
    \end{example}
    \begin{exercise}
        $\mathbb Z$ به‌عنوان $\mathbb R$-مدول؟
    \end{exercise}
\end{frame}

% Source: sources/lectures/04 Module.pdf | PDF page: 2 | item: commutative rings and restriction of scalars
\begin{frame}{تغییر سمت عمل و تغییر حلقهٔ اسکالر}
    \label{lec04:restriction-scalars}
    \begin{remark}
        اگر $R$ جابجایی باشد، هر $R$-مدول چپ را با تعریف
        \[
            m\cdot r:=rm
        \]
        می‌توان یک $R$-مدول راست در نظر گرفت.
        آیا شرط جابجایی لازم بود؟
    \end{remark}
    \begin{remark}
        اگر $\varphi:R\to S$ هم‌ریختی حلقه‌ای و $M$ یک $S$-مدول چپ باشد،
        با تعریف
        \[
            r\cdot m:=\varphi(r)m
        \]
        $M$ یک $R$-مدول چپ نیز هست.
    \end{remark}
\end{frame}

% Source: sources/lectures/04 Module.pdf | PDF page: 2-3 | item: submodule and trivial examples
\begin{frame}{زیرمدول}
    \label{lec04:submodule}
    \begin{definition}
        فرض کنید $M$ یک $R$-مدول چپ باشد.
        زیرگروه $(N,+)$ از $(M,+)$ را \textbf{زیرمدول} $M$ می‌نامیم هرگاه
        \[
            rn\in N\qquad(\forall r\in R,\ \forall n\in N).
        \]
    \end{definition}
    \begin{example}
        $\{0\}$ و $M$ زیرمدول‌های بدیهی $M$ هستند.
    \end{example}
\end{frame}

% Source: sources/lectures/04 Module.pdf | PDF page: 3 | item: quotient module and verification questions
\begin{frame}{مدول خارج‌قسمتی}
    \label{lec04:quotient-module}
    \begin{definition}
        اگر $N$ زیرمدول $M$ باشد، $N$ زیرگروهی نرمال از گروه آبلی $M$ است.
        همدسته‌های $M/N$ گروهی جمعی تشکیل می‌دهند. تعریف می‌کنیم:
        \[
            r(m+N):=rm+N.
        \]
        با این عمل، $M/N$ یک $R$-مدول چپ می‌شود.
    \end{definition}
    \begin{exercise}
        چرا این عمل، $M/N$ را به مدول تبدیل می‌کند؟
        آیا این ساخت برای هر زیرمدول صحیح است؟
    \end{exercise}
\end{frame}

% Source: sources/lectures/04 Module.pdf | PDF page: 3 | item: module homomorphisms and related types
\begin{frame}{هم‌ریختی مدول‌ها}
    \label{lec04:module-homomorphism}
    \begin{definition}
        برای دو $R$-مدول چپ $M_1$ و $M_2$، نگاشت
        $\varphi:M_1\to M_2$ یک \textbf{$R$-هم‌ریختی} است هرگاه
        \[
        \begin{aligned}
            \varphi(m+m')&=\varphi(m)+\varphi(m')
                &&(m,m'\in M_1),\\
            \varphi(rm)&=r\varphi(m)
                &&(r\in R,\ m\in M_1).
        \end{aligned}
        \]
        به همین نحو، $R$-تک‌ریختی، $R$-بروریختی و $R$-یکریختی
        تعریف می‌شوند.
    \end{definition}
    \begin{exercise}
        نشان دهید $\ker\varphi$ زیرمدول $M_1$ و
        $\operatorname{Im}\varphi$ زیرمدول $M_2$ است.
    \end{exercise}
\end{frame}

% Source: sources/lectures/04 Module.pdf | PDF page: 4 | item: first isomorphism theorem
\begin{frame}{قضیهٔ اول یکریختی مدول‌ها}
    \label{lec04:first-isomorphism}
    \begin{theorem}
        اگر $\varphi:M\to M'$ یک $R$-هم‌ریختی باشد، آنگاه
        \[
            \frac{M}{\ker\varphi}\modulecong\operatorname{Im}\varphi.
        \]
    \end{theorem}
\end{frame}

% Source: sources/lectures/04 Module.pdf | PDF page: 4 | item: second isomorphism theorem and the map given in its proof
\begin{frame}{قضیهٔ دوم یکریختی مدول‌ها}
    \label{lec04:second-isomorphism}
    \begin{theorem}
        اگر $N$ و $K$ دو زیرمدول از $R$-مدول $M$ باشند، آنگاه
        \[
            \frac{N+K}{N}\modulecong\frac{K}{N\cap K}.
        \]
    \end{theorem}
    \textbf{نگاشت مورد استفاده در اثبات.}
    \[
        \varphi:K\longrightarrow\frac{N+K}{N},
        \qquad k\longmapsto k+N.
    \]
\end{frame}

% Source: sources/lectures/04 Module.pdf | PDF page: 4 | item: intersection of submodules
\begin{frame}{اشتراک زیرمدول‌ها}
    \label{lec04:intersection-submodules}
    \begin{lemma}
        اگر $\{M_i\}_{i\in I}$ خانواده‌ای از زیرمدول‌های $R$-مدول $M$ باشد،
        آنگاه
        \[
            \bigcap_{i\in I}M_i
        \]
        نیز زیرمدول $M$ است.
    \end{lemma}
\end{frame}

% Source: sources/lectures/04 Module.pdf | PDF page: 5 | item: generated submodule, empty set, cyclic and finite generation
\begin{frame}{زیرمدول تولیدشده}
    \label{lec04:generated-submodule}
    \begin{definition}
        اگر $X\subseteq M$، کوچکترین زیرمدول $M$ که شامل $X$ باشد،
        \textbf{زیرمدول تولیدشده توسط $X$} نام دارد و با
        $\langle X\rangle$ نشان داده می‌شود.
        \[
            \langle\varnothing\rangle=\{0\},\qquad
            \langle\{a\}\rangle=\langle a\rangle.
        \]
        زیرمدول $\langle a\rangle$ را \textbf{دوری} می‌نامیم.
        اگر $X$ متناهی باشد، $\langle X\rangle$ با تولید متناهی است.
    \end{definition}
    \begin{lemma}
        \[
            \langle X\rangle=
            \bigcap_{\substack{N\leq M\\X\subseteq N}}N.
        \]
    \end{lemma}
\end{frame}

% Source: sources/lectures/04 Module.pdf | PDF page: 5 | item: cyclic examples in Z and Z_30
\begin{frame}{نمونه‌های زیرمدول دوری}
    \label{lec04:cyclic-examples}
    \begin{example}
        در $\mathbb Z$ به‌عنوان $\mathbb Z$-مدول،
        \[
            \langle n\rangle=n\mathbb Z.
        \]
    \end{example}
    \begin{example}
        در $\mathbb Z_{30}$ به‌عنوان $\mathbb Z$-مدول،
        \[
            \langle\overline5\rangle=
            \{\overline0,\overline5,\overline{10},\overline{15},\overline{20},\overline{25}\},
        \]
        \[
            \langle\overline6\rangle=
            \{\overline0,\overline6,\overline{12},\overline{18},\overline{24}\}.
        \]
    \end{example}
\end{frame}

% Source: sources/lectures/04 Module.pdf | PDF page: 5 | item: finite generation of submodules and cyclic submodules
\begin{frame}{زیرمدول با تولید متناهی}
    \label{lec04:finite-generation}
    \begin{remark}
        اگر $N\leq M$ و مجموعهٔ متناهی $X\subseteq N$ وجود داشته باشد که
        \[
            N=\langle X\rangle,
        \]
        آنگاه $N$ یک زیرمدول با تولید متناهی است.
        اگر $X$ تک‌عضوی باشد، $N$ دوری است و می‌نویسیم $N=\langle a\rangle$.
    \end{remark}
\end{frame}

% Source: sources/lectures/04 Module.pdf | PDF page: 5-6 | item: Ra and the formula for a generated cyclic submodule
\begin{frame}{مقایسهٔ دو زیرمدول}
    \label{lec04:ra-versus-generated}
    \begin{definition}
        فرض کنید $R$ حلقه‌ای جابجایی و یکدار و $M$ یک $R$-مدول باشد.
        برای $a\in M$، مجموعهٔ زیر یک زیرمدول $M$ است:
        \[
            Ra=\{ra\mid r\in R\}.
        \]
    \end{definition}
    \begin{remark}
        آیا همواره $\langle a\rangle=Ra$؟ در حالت کلی:
        \[
            \langle a\rangle=\{ra+na\mid r\in R,\ n\in\mathbb Z\}.
        \]
    \end{remark}
    % AMBIGUITY[L04-Q01]: primary, PDF pages 5-6; unitary action is an explicit provisional working assumption, not a resolved reading.
    \textbf{خوانش موقت:} اگر $M$ یکانی باشد، $\langle a\rangle=Ra$.
    فرض یکانی‌بودن در این عبارت منبع صریح نیست. \lr{[L04-Q01]}
\end{frame}

% Source: sources/lectures/04 Module.pdf | PDF page: 6 | item: nonunital example, convention and epimorphism
\begin{frame}{مثالی با حلقهٔ غیریکدار}
    \label{lec04:nonunital-example}
    \begin{example}
        حلقهٔ غیریکدار $R=(5\mathbb Z,+,\cdot)$ را در نظر بگیرید.
        $\mathbb Z$ با ضرب معمول یک $R$-مدول است. برای $a=3$ داریم:
        \[
            Ra=15\mathbb Z,\qquad \langle3\rangle=3\mathbb Z.
        \]
    \end{example}
    در ادامهٔ این بحث، حلقه‌ها جابجایی و یکدار فرض می‌شوند.
    \begin{remark}
        نگاشت زیر یک $R$-بروریختی است:
        \[
            \varphi:R\longrightarrow Ra,\qquad r\longmapsto ra.
        \]
    \end{remark}
\end{frame}

% Source: sources/lectures/04 Module.pdf | PDF page: 6 | item: structure theorem for cyclic unitary modules
\begin{frame}{ساختار مدول‌های دوری}
    \label{lec04:cyclic-structure}
    \begin{theorem}
        فرض کنید $R$ حلقه‌ای یکدار و $M$ یک $R$-مدول یکانی باشد.
        $M$ دوری است اگر و تنها اگر برای یک ایده‌آل چپ $I$ از $R$،
        \[
            M\modulecong R/I.
        \]
        اگر $M=\langle a\rangle$، می‌توان گرفت:
        \[
            I=\operatorname{Ann}_R(a)=\{r\in R\mid ra=0\}.
        \]
    \end{theorem}
\end{frame}

% Source: sources/lectures/04 Module.pdf | PDF page: 6-7 | item: forward proof of cyclic module structure theorem
% AMBIGUITY[L04-Q02]: page 6 bottom literally says a in R; use a in M provisionally and preserve the question in its PDF.
\begin{frame}{اثبات قضیهٔ ساختار: نگاشت طبیعی}
    \label{lec04:cyclic-proof-forward}
    \textbf{خوانش موقت:} مولد را عضو $M$ می‌گیریم؛ در منبع $a\in R$
    نوشته شده است. \lr{[L04-Q02]}

    \begin{proof}
        اگر $M$ دوری باشد، چون $R$ یکدار و $M$ یکانی است،
        برای یک $a\in M$ داریم $M=Ra$.
        نگاشت
        \[
            \varphi:R\longrightarrow M=Ra,\qquad r\longmapsto ra
        \]
        یک $R$-بروریختی است. بنابراین
        \[
            R/\ker\varphi\modulecong\operatorname{Im}\varphi=M.
        \]
        همچنین
        \[
            \ker\varphi=\operatorname{Ann}_R(a).
        \]
    \end{proof}
\end{frame}

% Source: sources/lectures/04 Module.pdf | PDF page: 7 | item: converse proof of cyclic module structure theorem
% AMBIGUITY[L04-Q03]: source writes phi(r) for an R/I-domain map; evaluations on cosets are provisional notation, not a recorded user resolution.
\begin{frame}{اثبات قضیهٔ ساختار: جهت عکس}
    \label{lec04:cyclic-proof-converse}
    \textbf{خوانش موقت:} در نمادگذاری نگاشت، $\varphi(r+I)$ به‌جای
    $\varphi(r)$ نوشته می‌شود. \lr{[L04-Q03]}

    \begin{proof}
        اگر $I$ ایده‌آلی چپ از $R$ و $M\modulecong R/I$ باشد،
        یک $R$-یکریختی $\varphi:R/I\to M$ در نظر می‌گیریم.
        با قرار دادن $a=\varphi(1+I)$ داریم:
        \[
            \varphi(r+I)=\varphi\bigl(r(1+I)\bigr)
            =r\varphi(1+I)=ra.
        \]
        پس $M$ توسط $\{a\}$ تولید می‌شود و دوری است.
    \end{proof}
\end{frame}

% Source: sources/lectures/04 Module.pdf | PDF page: 7 | item: sum of submodules and generated union
\begin{frame}{مجموع زیرمدول‌ها}
    \label{lec04:sum-submodules}
    \begin{definition}
        اگر $\{M_i\}_{i\in I}$ خانواده‌ای از زیرمدول‌های $R$-مدول $M$ باشد،
        مجموع آن‌ها را چنین تعریف می‌کنیم:
        \[
            \sum_{i\in I}M_i=
            \left\{a_1+\cdots+a_n\ \middle|\
            n\in\mathbb N,\ a_j\in M_{i_j}\right\}.
        \]
    \end{definition}
    \begin{lemma}
        \[
            \sum_{i\in I}M_i=
            \left\langle\bigcup_{i\in I}M_i\right\rangle.
        \]
    \end{lemma}
\end{frame}

% Source: sources/lectures/04 Module.pdf | PDF page: 8 | item: direct product and coordinatewise operations
\begin{frame}{حاصل‌ضرب مستقیم مدول‌ها}
    \label{lec04:direct-product}
    \begin{definition}
        فرض کنید $\{M_i\}_{i\in I}$ خانواده‌ای از $R$-مدول‌ها باشد.
        خانواده‌های $(m_i)_{i\in I}$ را در نظر می‌گیریم که $m_i\in M_i$.
        عملیات را مؤلفه‌به‌مؤلفه تعریف می‌کنیم:
        \[
            (m_i)_{i\in I}+(m_i')_{i\in I}
            =(m_i+m_i')_{i\in I},
        \]
        \[
            r(m_i)_{i\in I}=(rm_i)_{i\in I}.
        \]
        مدول حاصل، \textbf{حاصل‌ضرب مستقیم} مدول‌های $M_i$ نام دارد:
        \[
            \prod_{i\in I}M_i.
        \]
    \end{definition}
\end{frame}

% Source: sources/lectures/04 Module.pdf | PDF page: 8 | item: external direct sum and the finite index case
\begin{frame}{حاصل‌جمع مستقیم مدول‌ها}
    \label{lec04:external-direct-sum}
    \begin{definition}
        خانواده‌هایی از $\prod_{i\in I}M_i$ که فقط تعداد متناهی از
        مؤلفه‌های آن‌ها ناصفر است، یک $R$-مدول تشکیل می‌دهند.
        آن را \textbf{حاصل‌جمع مستقیم} می‌نامیم و با نماد
        \[
            \bigoplus_{i\in I}M_i
        \]
        نشان می‌دهیم.
    \end{definition}
    \[
        \bigoplus_{i\in I}M_i\leq\prod_{i\in I}M_i.
    \]
    اگر $I$ متناهی باشد،
    \[
        \bigoplus_{i\in I}M_i=\prod_{i\in I}M_i.
    \]
\end{frame}

% Source: sources/lectures/04 Module.pdf | PDF page: 9 | item: natural surjection from the external direct sum
\begin{frame}{نگاشت طبیعی از جمع مستقیم}
    \label{lec04:natural-sum-map}
    فرض کنید $\{M_i\}_{i\in I}$ خانواده‌ای از زیرمدول‌های $R$-مدول $M$ باشد.
    با تعریف
    \[
        \varphi:\bigoplus_{i\in I}M_i\longrightarrow\sum_{i\in I}M_i,
        \qquad (m_i)_{i\in I}\longmapsto\sum_{i\in I}m_i,
    \]
    یک $R$-هم‌ریختی پوشا به دست می‌آید. بنابراین
    \[
        \sum_{i\in I}M_i\modulecong
        \frac{\bigoplus_{i\in I}M_i}{\ker\varphi}.
    \]
\end{frame}

% Source: sources/lectures/04 Module.pdf | PDF page: 9 | item: trivial kernel and unique finite-sum representations
\begin{frame}{هسته و یکتایی نمایش}
    \label{lec04:kernel-uniqueness}
    برای نگاشت طبیعی $\varphi$، شرط $\ker\varphi=\{0\}$ معادل
    یکتایی نمایش عناصر مجموع است.
    به بیان دیگر، برای اندیس‌های متمایز $i_1,\ldots,i_k$ و $m_{i_j}\in M_{i_j}$،
    \[
        m_{i_1}+\cdots+m_{i_k}=0
        \quad\Longrightarrow\quad
        m_{i_j}=0\qquad(1\leq j\leq k).
    \]
    % AMBIGUITY[L04-Q04]: page 9 lowercase m_j is read as submodules M_j provisionally.
    \medskip
    \textbf{خوانش موقت:} با خواندن $M_j$ به‌جای $m_j$ در عبارت منبع،
    شرط بالا معادل است با \lr{[L04-Q04]}
    \[
        M_i\cap\sum_{\substack{j\in I\\j\ne i}}M_j=\{0\}
        \qquad(\forall i\in I).
    \]
\end{frame}

% Source: sources/lectures/04 Module.pdf | PDF page: 9-10 | item: internal direct sum definition
\begin{frame}{مجموع مستقیم داخلی}
    \label{lec04:internal-direct-sum}
    \begin{definition}[\lr{Internal direct sum}]
        فرض کنید $\{M_i\}_{i\in I}$ خانواده‌ای از زیرمدول‌های $R$-مدول $M$ باشد.
        $M$ را \textbf{مجموع مستقیم داخلی} این زیرمدول‌ها می‌نامیم هرگاه
        \[
            M=\sum_{i\in I}M_i
        \]
        و هر عنصر $M$، نمایش یکتایی به‌صورت مجموع متناهی از
        عناصر زیرمدول‌های $M_i$ داشته باشد.
        در این صورت می‌نویسیم:
        \[
            M=\bigoplus_{i\in I}M_i.
        \]
    \end{definition}
\end{frame}

% Source: sources/lectures/04 Module.pdf | PDF page: 10 | item: direct-sum criterion
% AMBIGUITY[L04-Q04]: the source omits spanning; the hypothesis M=sum M_i is explicitly provisional.
\begin{frame}{معیار مجموع مستقیم}
    \label{lec04:direct-sum-criterion}
    \textbf{فرض موقت افزوده‌شده:} $M=\sum_{i\in I}M_i$.
    این فرض در صورت نتیجهٔ دست‌نویس صریح نیست. \lr{[L04-Q04]}

    \begin{corollary}[صورت موقت]
        برای خانوادهٔ زیرمدول‌های $\{M_i\}_{i\in I}$ از $M$، با فرض بالا،
        دو گزارهٔ زیر معادل‌اند:
        \[
            M=\bigoplus_{i\in I}M_i,
        \]
        \[
            M_i\cap\sum_{\substack{j\in I\\j\ne i}}M_j=\{0\}
            \qquad(\forall i\in I).
        \]
    \end{corollary}
\end{frame}

% Source: sources/lectures/04 Module.pdf | PDF page: 10 | item: direct sum implies zero intersections
% AMBIGUITY[L04-Q04]: proof belongs to the explicitly provisional criterion.
\begin{frame}{اثبات معیار: اشتراک صفر}
    \label{lec04:direct-sum-proof-intersection}
    \textbf{اثباتِ صورت موقت.} \lr{[L04-Q04]}
    فرض کنید $M=\bigoplus_{i\in I}M_i$ و
    \[
        a\in M_i\cap\sum_{\substack{j\in I\\j\ne i}}M_j.
    \]
    در این صورت، برای اندیس‌هایی غیر از $i$،
    \[
        a=m_i=m_{t_1}+\cdots+m_{t_k}.
    \]
    بنابراین
    \[
        m_i-m_{t_1}-\cdots-m_{t_k}=0+0+\cdots+0.
    \]
    از یکتایی نمایش نتیجه می‌شود $a=m_i=0$.
    \hfill$\qed$
\end{frame}

% Source: sources/lectures/04 Module.pdf | PDF page: 10-11 | item: zero intersections imply unique representation
% AMBIGUITY[L04-Q04]: proof uses the provisional spanning hypothesis recorded with the criterion.
\begin{frame}{اثبات معیار: یکتایی نمایش}
    \label{lec04:direct-sum-proof-uniqueness}
    \textbf{ادامهٔ اثباتِ صورت موقت.} \lr{[L04-Q04]}
    برعکس، فرض کنید شرط اشتراک صفر برقرار باشد و یک عنصر دو نمایش داشته باشد:
    \[
        m_{t_1}+\cdots+m_{t_k}=m'_{s_1}+\cdots+m'_{s_r}.
    \]
    با یکی‌گرفتن اندیس‌ها داریم:
    \[
        m_{t_1}-m'_{t_1}
        =(m'_{t_2}-m_{t_2})+\cdots+(m'_{t_k}-m_{t_k}).
    \]
    طرف چپ عضو $M_{t_1}$ و طرف راست عضو
    $\sum_{j\ne t_1}M_j$ است؛ پس هر دو صفرند و
    $m_{t_1}=m'_{t_1}$. همین استدلال برای بقیهٔ مؤلفه‌ها برقرار است.
    در نتیجه نمایش یکتا است و با فرض تولید، $M=\bigoplus_{i\in I}M_i$.
    \hfill$\qed$
\end{frame}

% Source: sources/lectures/04 Module.pdf | PDF page: 11 | item: generating set as sum of cyclic submodules
\begin{frame}{مجموعهٔ مولد و مجموع زیرمدول‌های دوری}
    \label{lec04:generating-set-sum}
    \begin{remark}
        فرض کنید $R$ حلقه‌ای جابجایی و یکدار و $M$ یک $R$-مدول یکانی باشد.
        اگر $X$ مجموعه‌ای مولد برای $M$ باشد، آنگاه
        \[
            M=\sum_{x\in X}Rx.
        \]
    \end{remark}
\end{frame}

% Source: sources/lectures/04 Module.pdf | PDF page: 11 | item: maximal submodule definition and pZ example
\begin{frame}{زیرمدول ماکسیمال}
    \label{lec04:maximal-submodule}
    \begin{definition}
        زیرمدول $N$ از $R$-مدول $M$ را \textbf{ماکسیمال} می‌نامیم هرگاه
        $N\neq M$ و برای هر زیرمدول $K$ از $M$،
        \[
            N\subseteq K\subseteq M
            \quad\Longrightarrow\quad K=N\ \text{یا}\ K=M.
        \]
    \end{definition}
    \begin{example}
        در $\mathbb Z$ به‌عنوان $\mathbb Z$-مدول، اگر $p$ عددی اول باشد،
        $N=p\mathbb Z$ یک زیرمدول ماکسیمال است.
    \end{example}
\end{frame}

% Legacy introductory comments preserved verbatim; audited replacements appear above.


% \subsection*{}


% \begin{definition}[مدول]
%     فرض کنید $R$ یک حلقه باشد.
%     $(M,+)$ یک گروه آبلی باشد، همانند فضاهای برداری  یک ضرب اسکالر به‌صورت زیر تعریف می‌کنیم:
%     \[
%         \forall r \in R,\ \forall m \in M; \quad r.m \in M
%     \]

%     که در خواص زیر صدق کند:
%     $$\forall m_1, m_2 \in M,\ \forall r \in R;\quad r(m_1 + m_2) = rm_1 + rm_2$$
%     $$\forall r, s \in R,\ \forall m \in M;\quad (r+s)m = rm + sm$$
%     $$\forall r, s \in R,\ \forall m \in M;\quad (rs)m = r(sm)$$


%     در این صورت $M$ را یک
%     \textbf{$R$-مدول چپ}
%     می‌نامیم.
%     به‌صورت مشابه،
%     \textbf{$R$-مدول راست}
%     تعریف می‌شود.

% \end{definition}

% \begin{definition}[$R$-مدول چپ یکانی]
%     اگر حلقه $R$
%     یکدار باشد
%     و $M$،
%     یک
%     $R$-مدول چپ
%     باشد که
%     \[
%         \forall m \in M;\quad 1.m = m
%     \]
%     آنگاه $M$ را یک
%     $R$-مدول چپ یکانی
%     می‌گویند.

% \end{definition}


% \begin{example}
%     هر فضای برداری بر میدان $F$ یک $F$-مدول چپ یکانی است.
% \end{example}


% \begin{example}
%     هر حلقه دلخواه $R$، با ضرب:
%     \[
%         \forall r, s \in R:\quad r.s = rs
%     \]
%     یک $R$-مدول چپ می‌شود.
% \end{example}



% \begin{example}
%     هر گروه آبلی یک $\mathbb{Z}$-مدول چپ است.
% \end{example}


% \[
%     n \cdot a =
%     \begin{cases}
%         a + \cdots + a    & n > 0 \\
%         0                 & n = 0 \\
%         -(a + \cdots + a) & n < 0
%     \end{cases}
% \]



% \begin{example}
%     هر گروه آبلی $(M,+)$ را می‌توان با تعریف \(r \cdot m = 0\) ، به یک $R$-مدول تبدیل کرد.
% \end{example}

% \begin{example}
%     $\mathbb{R}[x]$ به‌عنوان یک $\mathbb{R}$-مدول چپ.
% \end{example}


% \begin{example}
%     $\mathbb{R}$ به‌عنوان $\mathbb{Z}$-مدول.
% \end{example}

% \begin{exercise}
%     $\mathbb{Z}$ به‌عنوان $\mathbb{R}$-مدول؟
% \end{exercise}


% \begin{remark}
%     اگر $R$ جابجایی باشد، هر $R$-مدول چپ را می‌توان به‌صورت $R$-مدول راست درنظر گرفت، با تعریف:
%     \[
%         m \cdot r := r \cdot m
%         \quad
%     \]
% \end{remark}
% (شرط جابجایی نیاز بود؟)


% \begin{remark}
%     اگر $\varphi: R \to S$ یک هم‌ریختی حلقه‌ای باشد و $M$ یک $S$-مدول چپ باشد، آنگاه $M$، یک $R$-مدول چپ هم هست، به‌صورت:
%     \[
%         r \cdot m := \varphi(r) \cdot m
%     \]
% \end{remark}


% \begin{definition}[زیرمدول]
%     $M$ یک $R$-مدول چپ است، زیرگروه $(N,+) \subseteq (M,+)$ را یک زیرمدول $M$ گوییم هرگاه:
%     \[
%         \forall r \in R,\ \forall n \in N;\quad r \cdot n \in N
%     \]
% \end{definition}

% \begin{example}
%     \(\{0\}\)
%     و
%     \(M\)
%     زیرمدول‌های بدیهی هستند
% \end{example}


% \begin{definition}[{مدول خارج‌قسمتی}]
%     اگر $N$ یک زیرمدول از $M$ باشد،
%     در این صورت
%     $N$
%     زیرگروه نرمال
%     $M$
%     است و لذا همدسته‌ها در
%     $M/N$
%     تشکیل گروه جمعی می‌دهند.
%     تعریف می‌کنیم:
%     \[
%         r \cdot (m + N) := r \cdot m + N
%     \]
%     در این صورت $M/N$ یک $R$-مدول چپ می‌شود. (چرا؟)
% \end{definition}

% آیا برای هر زیرمدول صحیح است؟


% \begin{definition}[{$R$-هم‌ریختی}]
%     اگر $M_1$ و $M_2$ دو $R$-مدول چپ باشند، آنگاه نگاشت
%     \[
%         \varphi : M_1 \longrightarrow M_2
%     \]
%     یک \textbf{$R$-هم‌ریختی} نامیده می‌شود هرگاه:
%     \begin{align*}
%         \text{(1)}\quad & \forall m, m' \in M_1;\ \varphi(m + m') = \varphi(m) + \varphi(m')            \\
%         \text{(2)}\quad & \forall m \in M_1,\ \forall r \in R;\ \varphi(r \cdot m) = r \cdot \varphi(m)
%     \end{align*}
%     به همین نحو
%     \textbf{$R$-تک‌ریختی}،
%     \textbf{$R$-بروریختی}و
%     \textbf{$R$-یکریختی}
%     تعریف می‌شود.

% \end{definition}


% \begin{theorem}
%     اگر $\varphi : M \to M'$ یک $R$-هم‌ریختی باشد، آنگاه:
%     \begin{itemize}
%         \item $ker \varphi$ یک زیرمدول از $M$ است.
%         \item $Im  \varphi$ یک زیرمدول از $M'$ است.
%     \end{itemize}

% \end{theorem}


% \begin{theorem}[{قضیه اول هم‌ریختی مدول‌ها}]
%     اگر $\varphi : M \to M'$ یک $R$-هم‌ریختی باشد، آنگاه:
%     \[
%         \frac{M}{ker \varphi} \modulecong Im  \varphi
%     \]
% \end{theorem}


% \begin{theorem}[{قضیه دوم هم‌ریختی مدول‌ها}]
%     اگر $N$ و $K$ دو زیرمدول از $R$-مدول $M$ باشند، در این‌صورت:
%     \[
%         \frac{N + K}{N} \modulecong \frac{K}{N \cap K}
%     \]


% \end{theorem}

% \begin{proof}
%     {تعریف می‌کنیم:}
%     \[\varphi : K \to \frac{N + K}{N}, \quad k \mapsto k + N\]

% \end{proof}

% \begin{lemma}
%     اگر $\{ M_i \}_{i \in I}$ خانواده‌ای از زیرمدول‌های $R$-مدول $M$ باشند، آنگاه
%     \(\bigcap_{i \in I} M_i\)
%     نیز زیرمدول $M$ است.
% \end{lemma}



% Legacy active draft preserved as comments for comparison, not source-verified content.
% \begin{frame}{زیرمدول تولیدشده}
%     \begin{definition}
%         اگر $M$ یک $R$-مدول باشد و $X \subseteq M$ آنگاه کوچکترین $R$-زیرمدول $M$ که شامل $X$ باشد را زیرمدول تولیدشده توسط $X$ می‌نامند و آن را با $\langle X \rangle$ نشان می‌دهند.
%         اگر $X = \{ a \}$، می‌نویسیم $\langle a \rangle$.
%
%         اگر
%         \(X = \emptyset\)
%         به‌وضوح
%         \(\langle X \rangle = \{0\}\).
%         اگر
%         \(X = {a}\)
%         آنگاه
%         \(\langle a \rangle\)
%         یک زیرمدول دوری
%         \(M\)
%         نامیده می‌شود. اگر
%         \(X\)
%         متناهی باشد
%         $\langle X \rangle$
%         با تولید متناهی نامیده می‌شود.
%     \end{definition}
%
%     \begin{lemma}
%         با فرضیات فوق داریم:
%         \[
%             \langle X \rangle = \bigcap_{\substack{N \leqslant M ,  X \subseteq N}} N
%         \]
%     \end{lemma}
%
%
% \end{frame}
%
%
% \begin{frame}{نمونه‌های زیرمدول دوری}
%     \begin{example}
%         در $\mathbb{Z}$ به‌عنوان $\mathbb{Z}$-مدول:
%         \[
%             \langle n \rangle = n \mathbb{Z}
%         \]
%
%     \end{example}
%
%     \begin{example}
%         در $\mathbb{Z}_{30}$ به‌عنوان $\mathbb{Z}$-مدول:
%         \[
%             \langle \overline{5} \rangle = \{ \overline{0}, \overline{5}, \overline{10}, \overline{15}, \overline{20}, \overline{25} \}
%         \]
%         \[
%             \langle \overline{4} \rangle = \{ \overline{0}, \overline{4}, \overline{8}, \overline{12}, \overline{16}, \overline{20}, \overline{24}, \overline{28} \}
%         \]
%     \end{example}
%
%
% \end{frame}
%
% \begin{frame}{زیرمدول‌های دوری}
%
%     \begin{definition}
%         اگر $R$ حلقه جابجایی و یکدار و $M$ یک $R$-مدول باشد، برای $a \in M$ تعریف می‌کنیم:
%         \[
%             Ra = \{ ra \mid r \in R \}
%         \]
%         که یک زیرمدول از $M$ خواهد بود.
%
%     \end{definition}
%
% \end{frame}
%
% \begin{frame}{نقش عضو همانی در تولید زیرمدول}
%     \begin{remark}
%         آیا $\langle a \rangle = Ra$ همواره برقرار است؟
%     \end{remark}
%
%
%     \[
%         Ra = \{ ra \mid r \in R \}
%         \qquad ,
%         \langle a \rangle = \{ ra + na \mid r \in R,\ n \in \mathbb{Z} \}
%     \]
%
%     ولی اگر حلقه
%     یکدار باشد،‌ یکی هستند.
%     به عنوان مثال حلقهٔ غیر یکدار
%     \(R = (5\mathbb{Z}, +, \cdot)\)
%     را در نظر بگیرید.
%     به وضوح
%     \(\mathbb{Z}\)
%     یک
%     \(R\)-مدول
%     است. حالا اگر
%     \(a = 3\)
%     \[
%         Ra = \{ ra \mid r \in R \} = 15\mathbb{Z}
%     \]
%     در حالی که
%     \(\langle 3 \rangle = 3 \mathbb{Z}\). البته در ادامه همواره حلقه های را جابجایی و یکدار فرض می‌کنیم.
%
%     \begin{remark}
%         با
%         تعریف
%         \[
%             \varphi : R \to Ra, \quad r \mapsto ra
%         \]
%         به یک
%         $R$-بروریختی
%         می‌رسیم.
%
%     \end{remark}
%
% \end{frame}
%
%
%
% \begin{frame}{ساختار مدول‌های دوری}
%     \begin{theorem}
%         فرض کنید $R$ حلقه‌ای یکدار باشد. در این صورت،
%         \(R\)-مدول
%         یکانی
%         \(M\)
%         دوری است اگر و تنها اگر
%         \(M \cong \frac{R}{I}\)
%         که در آن
%         \(I\)
%         یک ایده‌آل چپ
%         \(R\)
%         است.
%         [در واقع \(I = Ann_R(a) = \{ r \in R \mid ra = 0 \}\)]
%     \end{theorem}
%
%     \[
%         M = \langle a \rangle \Rightarrow \frac{R}{Ann_R(a)} \modulecong M = Ra
%     \]
%
% \end{frame}
%
%
% \begin{frame}{اثبات قضیهٔ ساختار: نگاشت طبیعی}
%     \textbf{اثبات.}
%     چون $R$ یکدار است و $M$ یکانی است، پس
%     \[
%         \exists a \in R,\quad M = Ra
%     \]
%
%     در این صورت
%     \[
%         \varphi : R \longrightarrow M = Ra, \quad r \longmapsto ra
%     \]
%     $\varphi$ یک $R$-بروریختی  است. بنابراین:
%     \[
%         \frac{R}{ker \varphi} \modulecong Im  \varphi = M
%     \]
%     و واضح است که
%     \[
%         ker \varphi = Ann_R(a)
%     \]
%
%
% \end{frame}
%
% \begin{frame}{اثبات قضیهٔ ساختار: جهت عکس}
%
%     \textbf{برعکس.}
%     اگر $M \cong \frac{R}{I}$ که $I$ یک ایده‌آل چپ از $R$ باشد، فرض کنید:
%
%     \[
%         \varphi : \frac{R}{I} \longrightarrow M, \quad \varphi(r + I) = \varphi(r)
%     \]
%     $R$-یکریختی
%     موردنظر باشد.
%     در این صورت با فرض $a = \varphi(1 + I)$ داریم
%     \[
%         \varphi(r + I) = \varphi(r(1 + I)) = r \varphi(1 + I) = ra
%     \]
%     و لذا $M$ توسط $\{a\}$ تولید می‌شود.
%
% \end{frame}
%
%
% \begin{frame}{مجموع زیرمدول‌ها}
%     \begin{definition}
%         اگر $M$ یک $R$-مدول و $\{M_i\}_{i \in I}$ زیرمدول‌های آن باشند، مجموع زیرمدول‌ها را به‌صورت زیر تعریف می‌کنیم:
%
%         \[
%             \sum_{i \in I} M_i := \left\{ a_1 + a_2 + \cdots + a_n \mid n \in \mathbb{N},\ a_j \in M_{i_j} \right\}
%         \]
%
%
%
%     \end{definition}
%
%     \begin{lemma}
%         واضح است که
%         \[
%             \sum_{i \in I} M_i = \left\langle \bigcup_{i \in I} M_i \right\rangle
%         \]
%     \end{lemma}
%
% \end{frame}
%
% \begin{frame}{حاصل‌ضرب مستقیم مدول‌ها}
%     \begin{definition}
%         فرض کنید $\{M_i\}_{i \in I}$ خانواده‌ای از $R$-مدول‌ها باشند.
%         در این صورت دنباله‌هایی به شکل $\{m_i\}_{i \in I}$ را در نظر می‌گیریم که در آن $m_i \in M_i$. سپس عملیات‌های زیر را تعریف می‌کنیم:
%
%         \[
%             \{m_i\}_{i \in I} + \{m_i'\}_{i \in I} := \{m_i + m_i'\}_{i \in I}
%             \quad ; \quad
%             r \cdot \{m_i\}_{i \in I} = \{r m_i\}_{i \in I}
%         \]
%         با این جمع و ضرب تعریف‌شده، یک $R$-مدول بدست می‌آید که آن را حاصل‌ضرب مستقیم $M_i$ها می‌خوانند و با نماد
%         \[
%             \prod_{i \in I} M_i
%         \]
%         نمایش می‌دهیم.
%
%     \end{definition}
%
% \end{frame}
%
%
%
% \begin{frame}{حاصل‌جمع مستقیم مدول‌ها}
%     حال اگر فقط دنباله‌هایی به شکل
%     $\{M_i\}_{i \in I}$
%     را در نظر بگیریم که فقط تعداد متناهی از عناصر آن‌ها
%     ناصفرند،
%     به زیرمجموعه‌ای از قبلی می‌رسیم که خود یک
%     $R$-مدول
%     است و آن را
%     حاصل‌جمع مستقیم
%     $M_i$ ها
%     گویند و با علامت
%     \(\bigoplus_{i \in I} M_i\)
%     نمایش داده می‌شود.
%
%     واضح است که
%     $\bigoplus_{i \in I} M_i$
%     زیرمدولی از
%     $\prod_{i \in I} M_i$
%     است.
%     اگر $I$ متناهی باشد:
%     \[
%         \prod_{i \in I} M_i = \bigoplus_{i \in I} M_i.
%     \]
%
% \end{frame}
%
% \begin{frame}{نگاشت طبیعی از جمع مستقیم}
%     \begin{remark}
%         حال اگر فرض کنیم $\{M_i\}_{i \in I}$ خانواده‌ای از زیرمدول‌های یک $R$-مدول $M$ باشند، در این صورت با تعریف
%         \[
%             \varphi : \bigoplus_{i \in I} M_i \longrightarrow \sum_{i \in I} M_i \; , \; \{m_i\}_{i \in I} \longmapsto \sum_{i \in I} m_i
%         \]
%         \[
%         \]
%         یک $R$-همریختی پوشا می‌رسیم و لذا:
%         \[
%             \sum_{i \in I} M_i \ \cong_{\text{همریختی}-R} \ \frac{\bigoplus_{i \in I} M_i}{ker \varphi}.
%         \]
%
%     \end{remark}
% \end{frame}
%
% \begin{frame}{هستهٔ نگاشت و یکتایی نمایش}
%     اما در این صورت $ker \varphi = \{0\}$ به چه معنا است؟
%     \[
%         \begin{aligned}
%             ker \varphi = \{0\} \ \iff & \text{نمایش  منحصر بفرد}                                                            \\
%             \iff                        & \left( m_{i_1} + \cdots + m_{i_k} = 0 \ \Rightarrow \ \forall j : m_{i} = 0 \right) \\
%             \iff                        & \quad M_i \ \cap \ \sum_{\substack{j \in I                                          \\ j \neq i}} m_j \ = \ \{0\}.
%         \end{aligned}
%     \]
%
% \end{frame}
%
% \begin{frame}{مجموع مستقیم داخلی}
%
%     \begin{theorem}[\lr{Internal direct sum}]
%         فرض کنید $M$ یک $R$-مدول است و $\{M_i\}_{i \in I}$ خانواده‌ای از زیرمدول‌های $M$ باشند. در این صورت $M$ را \textbf{مجموع مستقیم} زیرمدول‌های $M_i$ می‌نامند هرگاه:
%         \[
%             M = \sum_{i \in I} M_i
%         \]
%         و هر عنصر $M$ را بتوان به‌صورت
%         منحصر بفردی به شکل مجموع عناصری از
%         \(M_i\)
%         نوشت.
%         در این صورت می‌نویسیم
%         \[
%             M = \bigoplus_{i \in I} M_i.
%         \]
%
%     \end{theorem}
%
% \end{frame}
%
%
% \begin{frame}{معیار مجموع مستقیم}
%     \begin{corollary}
%         اگر $\{M_i\}_{i \in I}$ خانواده‌ای از $R$-زیرمدول‌های $M$ باشند، گزاره‌های زیر معادلند:
%         \begin{enumerate}
%             \item[(1)]
%                   \[
%                       M = \bigoplus_{i \in I} M_i
%                   \]
%             \item[(2)]
%                   \[
%                       M_i \ \cap \ \sum_{\substack{j \in I \\ j \neq i}} M_j \ = \ \{0\} \quad (\forall i \in I)
%                   \]
%         \end{enumerate}
%     \end{corollary}
%
%     \textbf{اثبات.}
%     اگر
%     \[
%         a \in M_i \ \cap \ \sum_{\substack{j \in I \\ j \neq i}} M_j
%     \]
%     پس
%     \[
%         a = m_i = m_{t_1} + m_{t_2} + \cdots + m_{t_k}
%     \]
% \end{frame}
%
% \begin{frame}{یکتایی نمایش در مجموع مستقیم}
%
%     در این صورت:
%     \[
%         m_i - m_{t_1} - m_{t_2} - \cdots - m_{t_k} = 0 + 0 + \cdots + 0
%     \]
%     و با توجه به نمایش یکتا
%     $$a = m_i = 0$$
%
%     برعکس، اگر نمایش یک عنصر به‌ دو صورت توشته شود:
%     \[
%         m_{t_1} + m_{t_2} + \cdots + m_{t_k} \ = \ m'_{s_1} + m'_{s_2} + \cdots + m'_{s_r}
%     \]
%     با یکی گرفتن این اندیس‌ها نتیجه می‌گیریم که
%     \[
%         m_{t_1} - m'_{t_1}
%         = (m'_{t_2} - m_{t_2}) + \cdots + (m'_{t_k} - m_{t_k})
%     \]
%     که
%     \[
%         m_{t_1} - m'_{t_1} \in M_{t_1},
%         \quad
%         (m'_{t_2} - m_{t_2}) + \cdots + (m'_{t_k} - m_{t_k})
%         \in \sum_{\substack{j \in I \\ j \neq t_1}} M_j
%     \]
%     لذا صفر می‌شود و پس $m_{t_1} = m'_{t_1}$ و همین‌طور برای بقیهٔ اعضا.
%
% \end{frame}
%
% \begin{frame}{مجموعهٔ مولد و زیرمدول ماکسیمال}
%     \begin{remark}
%         فرض کنید
%         \(R\)
%         حلقه‌ای جابجایی و یک‌دار است.
%         اگر $M$ یک $R$-مدول و $X$ یک مجموعهٔ مولد $M$ باشد آنگاه خواهیم داشت:
%         \[
%             M = \sum_{x \in X} Rx.
%         \]
%
%     \end{remark}
%
%
%     \begin{definition}
%         زیرمدول $N$ از $R$-مدول $M$ را \textbf{ماکسیمال} گوییم هرگاه $N \neq M$ و اگر
%         \[
%             N \subseteq K \subseteq M
%         \]
%         که $K$ یک $R$-زیرمدول $M$ است، آنگاه
%         \[
%             K = N \quad \text{یا} \quad K = M.
%         \]
%
%     \end{definition}
%
%
%
% \end{frame}
%
% \begin{frame}{نمونه‌ای از زیرمدول ماکسیمال}
%     \begin{example}
%         $\mathbb{Z}$ به عنوان $\mathbb{Z}$-مدول و
%         \[
%             N = p\mathbb{Z}
%         \]
%         که $p$ عددی اول است.
%
%     \end{example}
% \end{frame}
